%======================================================================
\subsection{BW-PDM-02 --- The spec as built (24 Oct 2026, approved 24 Oct)}\label{sec:pdm02-spec}
%======================================================================

The 19 Oct draft of \eref{PDM-02} was written before the folder vault (\S\ref{sec:pdm-decision}), the vault client
(\S\ref{sec:pdm04-plan}) and review and release (\S\ref{sec:pdm05-plan}) were built. Since then the code gained the
journal fields of \eref{PDM-05} (release checklist and note, rejection reason, ECR and ECO records and numbers,
retire with a reason), the working folder (\file{work/<actor>/}, \file{.rev/}), Vault Settings and Move Vault. This
section restates the spec from the code. Where the code and the 19 Oct text disagreed, the documents now say what
the code does, and the difference is listed under \emph{For approval} below. No behaviour was added for the spec.

\paragraph{What the spec is.} Three files and their tests:
\begin{itemize}[nosep]
  \item \file{bruceworx/pdm.py}: the rules (revisions, lifecycle, change control, BOM); storage-free.
  \item \file{bruceworx/pdm\_folder.py}: files, locks and the journal; every operation is one journal event, replayed
        through \file{pdm.py}.
  \item \file{docs/pdm-openapi.yaml}, version 0.2.0-draft: 21 operations. Each state-changing operation's
        \code{operationId} is its journal operation, so a later server (Ops) and the folder vault name the same things.
        The Ops numbering call of \eref{PDM-03} (\code{POST /part-numbers}) is in it too.
  \item Tests: the 19 Oct contract \file{tests/test\_pdm.py}, also re-run on a folder, and the new
        \file{tests/test\_pdm\_contract.py} (48 tests): every row of table~\ref{tab:pdm-transitions} on the
        reference model and on a folder vault (with a reopen after each test, so the journal replays to the same
        states), every refusal in the right-hand column, and a check that the table, the OpenAPI operations and the
        journal operations are the same list.
\end{itemize}

{\small
\begin{longtable}{@{}l>{\raggedright\arraybackslash}p{3.2cm}>{\raggedright\arraybackslash}p{4.6cm}>{\raggedright\arraybackslash}p{5.6cm}@{}}
\caption{Every transition of the PDM spec. IDs are the test names in \file{test\_pdm\_contract.py}.}\label{tab:pdm-transitions}\\
\toprule
ID & Operation & From $\rightarrow$ to & Refused when\\
\midrule
\endfirsthead
\toprule
ID & Operation & From $\rightarrow$ to & Refused when\\
\midrule
\endhead
T1 & \code{create\_item} \newline \code{POST /items} & new $\rightarrow$ rev 01, WIP & the P/N exists; a drawing names no
  existing part\\
T2 & \code{submit} & WIP $\rightarrow$ In Review & not WIP; the file is checked out (folder)\\
T3 & \code{reject} & In Review (01, 02, \ldots) $\rightarrow$ next number in WIP; the rejected one is superseded &
  not In Review\\
T4 & \code{reject} & In Review (ECO letter) $\rightarrow$ WIP, same letter & not In Review\\
T5 & \code{approve} & In Review (pre-release) $\rightarrow$ Released as A; its file is frozen & not In Review; an ECO
  revision; already released\\
T6 & \code{raise\_ecr} \newline \code{POST /ecrs} & released items $\rightarrow$ ECR Open & an item not released, or
  obsolete; empty problem, no items, number taken (folder)\\
T7, T8 & \code{decide\_ecr} & ECR Open $\rightarrow$ Approved, or Rejected & already decided\\
T9 & \code{open\_eco} \newline \code{POST /ecos} & released + approved ECRs $\rightarrow$ next letter in WIP for every
  item and its drawings (retired drawings stay retired) & an ECR not Approved; an item obsolete or with a revision in
  WIP or In Review. A refusal changes nothing.\\
T10 & \code{submit} & WIP (ECO letter) $\rightarrow$ In Review & as T2\\
T11 & \code{release\_eco} & every ECO revision In Review $\rightarrow$ all Released together; older letters
  superseded; files frozen; ECO Released & the ECO is not Open; an item is not In Review\\
T12 & \code{retire} & Released $\rightarrow$ Obsolete; never changes again & not released; already obsolete; its next
  letter is in work under an ECO; a released assembly uses it; no reason (folder). The clients also refuse while a
  released item's file links to it.\\
F1 & \code{checkout} & WIP, unlocked $\rightarrow$ locked by me & not WIP (so never a released file); locked by
  someone else\\
F2 & \code{checkin} & locked by me $\rightarrow$ new version, unlocked (or kept) & not locked by me; not WIP. The
  same content adds no version and only unlocks. A comment is required (client).\\
F3 & \code{undo\_checkout} & locked by me $\rightarrow$ unlocked & ---\\
F4 & \code{undo\_checkout} \newline (\code{force}) & locked by someone else $\rightarrow$ unlocked, recorded &
  without \code{force}\\
B1 & \code{set\_bom} \newline \code{PUT /items/\{pn\}/bom} & BOM replaced & unknown child; quantity below 1; the BOM
  would contain itself (the old BOM is kept)\\
\bottomrule
\end{longtable}}

Figure~\ref{fig:lifecycle-sm} now shows T4 (an ECO letter rejected keeps its letter) and labels T11 \emph{release
ECO}: an ECO letter is never released by \emph{approve}.

\paragraph{Decision records.} Every decision is recorded on the revision (or the ECR or ECO) with who (\code{user@mac})
and when: \emph{approve} with the release note, the checklist as ticked (\file{release\_check.py}) and attachments;
\emph{reject} with its reason; \emph{retire} with its reason; an ECR with its problem, attachments and the decision's
reason; an ECO with its release note and one checklist per item. Attachments are write-once objects like file
versions. ECRs and ECOs are numbered per vault (\code{ECR-0001}, \code{ECO-0001}) when no number is given. Fields are
written only when set, so journals from before 23 Oct replay unchanged.

\paragraph{The vault folder as built.}
\begin{center}
\small
\begin{tabularx}{\linewidth}{@{}lX@{}}
\toprule
Path & Holds\\
\midrule
\file{<name>.bwvault} & The ref file: format, version, vault id, name. Written once.\\
\file{objects/<ab>/<sha256>.<ext>} & File versions and attachments; write-once, verified on every read.\\
\file{journal/<actor>.jsonl} & Append-only events, one writer per file; sync-service duplicates are merged by event id.\\
\file{work/<actor>/<P/N>.FCStd} & Working files (\S\ref{sec:pdm04-plan}); writable only while checked out to you.\\
\file{work/<actor>/.bw/state.json} & Per file: the version it started from.\\
\file{work/<actor>/.rev/<P/N>-<rev>-<sha8>.FCStd} & Older versions from \emph{Get revision}, read-only, links pointed
  at the work copies.\\
\file{numbering.json}, \file{document\_types.json} & Part numbering (\S\ref{sec:pdm03}); Vault Settings adds, renames
  and retires rows, never deletes them or changes a code.\\
\file{part\_properties.json}, \file{drawing\_standards.json} & Company part properties; title-block standards.\\
\bottomrule
\end{tabularx}
\end{center}
\emph{Move Vault} copies the whole folder, compares every file by SHA-256, opens the copy and requires it to replay to
the same items, versions and locks, and only then renames it into place (staged as \file{<name>.moving}); the old
folder is left as it was. None of this changes the API: the work folder, settings files and moving are the folder
vault's business, not operations of the spec.

\paragraph{Your answers (24 Oct).} A1--A11 all as recommended. Approved as built: A1 (the spec is the three files
and table~\ref{tab:pdm-transitions}), A2, A4 (per-revision BOMs come with \eref{PDM-06}), A6, A7, A8 and A11. To
change: \textbf{A3}, the required texts (check-in comment, reject reason, release note, ECR problem, decision reason,
retire reason) move into \file{pdm.py} and the folder's check-in, so every client gets the same rule; \textbf{A9}, the
folder-only refusals that need no files (empty problem, no items, number taken, no retire reason) move into
\file{pdm.py} too; \textbf{A5}, the ECO state \emph{Cancelled} is dropped from the spec until a \emph{Cancel ECO} is
designed; \textbf{A10}, the file-link check of ``used by'' (\file{release\_check.where\_used}) becomes part of the spec,
until \eref{PDM-06} pushes BOMs on check-in and the BOM check alone is enough.

\paragraph{For approval (approved 24 Oct).} What you approve, with my recommendation.
\begin{center}
\small
\begin{tabularx}{\linewidth}{@{}l>{\raggedright\arraybackslash}p{6.2cm}X@{}}
\toprule
\# & Item & Recommendation\\
\midrule
A1 & The spec is \file{pdm.py} + \file{pdm\_folder.py} + \file{pdm-openapi.yaml} 0.2.0-draft, with
  table~\ref{tab:pdm-transitions} as its state machine. & Approve. It is what runs today.\\
A2 & Locks never expire; anyone can break one with \code{force}, and it is recorded. The draft had leases and an
  administrator. & Approve (your 23 Oct decision). Revisit roles with several users.\\
A3 & Required texts: the clients (FreeCAD panels; \code{bw-hooks vault approve}, \code{reject}, \code{retire
  --note}) require a check-in comment, a reject reason, a release note, an ECR problem, a decision reason and a retire
  reason; the folder vault enforces only the problem and the retire reason; the
  reference model none. & Move all of them into \file{pdm.py} (and the folder's check-in), so every client gets the
  same rule.\\
A4 & One BOM per item, not per revision; find number, configuration and pinned revision are not stored. The figure
  draws BOM lines from a parent \emph{revision}. & Approve for now; per-revision BOMs and those fields belong to
  \eref{PDM-06}.\\
A5 & The draft's ECO state \emph{Cancelled} has no operation. & Drop it from the spec (done in the YAML) until a
  \emph{Cancel ECO} is designed: what happens to its WIP letters is your call.\\
A6 & Retire as built: reason required, refused while the next letter is in work under an ECO, an obsolete item never
  changes again, a part's drawing is not retired with it (the dialog offers it). & Approve.\\
A7 & Unreleased work cannot be retired or deleted: it ``simply stops'' and stays listed in WIP. & Approve; add a
  \emph{Cancel} only if the WIP list gets cluttered.\\
A8 & Submit is refused while the file is checked out (folder). & Approve; the reviewed file must be the vault's.\\
A9 & Refusals that live only in the folder vault (empty problem, no items, number taken, no retire reason, submit while
  checked out) are not in the reference model, so a server written from \file{pdm.py} alone would accept them. &
  Move the ones that need no files into \file{pdm.py}.\\
A10 & ``Used by'' for retiring: the engine looks at BOM lines only; the clients also count released items whose
  checked-in file links to the part (\file{release\_check.where\_used}), so the API and FreeCAD can disagree. &
  Make the file-link check part of the spec: \eref{PDM-06} pushes BOMs on check-in, after which the BOM check alone
  is enough.\\
A11 & \emph{Get revision} names files \file{<P/N>-<rev>-<sha8>.FCStd}, not \file{<P/N>-<rev>.FCStd} as the
  \eref{PDM-04} plan says, so two versions of one revision can sit side by side. & Approve; the plan text follows.\\
\bottomrule
\end{tabularx}
\end{center}
